\chapter{The Singular Value Decomposition}

Here is the geometric fact at the center of this chapter:

\begin{quote}
\emph{Every matrix -- any matrix at all, square or rectangular, full rank or
rank one, invertible or hopelessly singular -- can be written as a rotation,
then a stretch along perpendicular axes, then another rotation.}
\end{quote}

That is the singular value decomposition, the SVD, and it is the most widely applied
factorization in numerical linear algebra.

In the previous chapter, we examined similarity and the Jordan canonical
form. There, we diagonalized a matrix using the \emph{same} basis for
both the input and output spaces: $\mat{P}\inv\mat{A}\mat{P}$. That demand was
restrictive. When a matrix sheared space, independent eigenvectors were missing,
forcing us to settle for Jordan blocks with off-diagonal ones. Even when
eigenvectors existed, they were often non-orthogonal.
Eigenvalues are powerful, but they have inherent limitations: they apply only to square matrices,
they can become complex even for real matrices, and defective matrices lack a full basis of eigenvectors.

The SVD avoids all of these restrictions. By choosing \emph{two different} orthonormal bases -- one for the input
space and one for the output space -- every matrix admits a clean diagonal representation.
It applies to any $m\times n$ matrix: a $3\times 7$ matrix of pixel intensities, a
$10000 \times 50$ matrix of customer ratings, or a $2 \times 2$
shear that has only one eigenvector. The decomposition always exists, its singular values are always
non-negative real numbers, and its geometric action is always the same: \textbf{rotate, stretch,
rotate}.

The SVD brings together the four fundamental subspaces, orthonormal bases,
projections, and the spectral theorem for symmetric matrices. When you finish
this chapter you will understand principal component analysis, low-rank
compression, the pseudoinverse, and a surprising amount of the linear algebra
that runs underneath modern machine learning -- all from one picture.

\section{The Picture First: What a Matrix Does to a Circle}
\label{ch29:circle-picture}

Forget formulas for a moment and let us just \emph{watch} a matrix act. Take

\[
  \mat{A} = \begin{bmatrix} 2 & 1 \\ 0 & 2 \end{bmatrix}.
\]

Feed it every unit vector in the plane -- the whole unit circle -- and look at
where the outputs land. A matrix sends straight lines to straight lines and the
unit circle (which is \emph{not} straight) to an ellipse. Always an ellipse.
Figure~\ref{ch29:fig-circle-to-ellipse} shows it.

\begin{figure}[H]
\centering
\begin{tikzpicture}[scale=0.9,>=Stealth]
  % left: unit circle. For A=[[2,1],[0,2]] the right singular vectors are
  % v1 at ~52deg and v2 at ~142deg (orthogonal).
  \begin{scope}[xshift=0cm]
    \draw[->,gray!60] (-1.6,0) -- (1.8,0) node[right]{$x$};
    \draw[->,gray!60] (0,-1.6) -- (0,1.8) node[above]{$y$};
    \draw[vecblue,thick] (0,0) circle (1);
    % two special input directions (the right singular vectors)
    \draw[->,vecred,very thick] (0,0) -- (0.615,0.788);
    \draw[->,structurepurple,very thick] (0,0) -- (-0.788,0.615);
    \node[vecred] at (0.95,1.0) {$\vec{v}_1$};
    \node[structurepurple] at (-1.1,0.75) {$\vec{v}_2$};
    \node at (0,-2.1) {input: unit circle};
  \end{scope}
  % arrow
  \draw[->,thick] (2.3,0) -- (3.5,0) node[midway,above]{$\mat{A}$};
  % right: ellipse
  \begin{scope}[xshift=7cm]
    \draw[->,gray!60] (-2.6,0) -- (3.4,0) node[right]{$x$};
    \draw[->,gray!60] (0,-2.6) -- (0,2.8) node[above]{$y$};
    % ellipse: A maps unit circle to ellipse with semi-axes sigma1=2.56, sigma2=1.56
    % the long axis u1 points at ~38deg.
    \draw[targetgreen,thick,rotate around={38:(0,0)}] (0,0) ellipse (2.56 and 1.56);
    \draw[->,vecred,very thick] (0,0) -- ({2.56*cos(38)},{2.56*sin(38)})
      node[above right=1pt, vecred, font=\small] {$\sigma_1 \vec{u}_1$};
    \draw[->,structurepurple,very thick] (0,0) -- ({-1.56*sin(38)},{1.56*cos(38)})
      node[above left=1pt, structurepurple, font=\small] {$\sigma_2 \vec{u}_2$};
    \node at (0,-2.9) {output: an ellipse};
  \end{scope}
\end{tikzpicture}
\caption{Any $2\times 2$ matrix sends the unit circle to an ellipse. The two
perpendicular input directions (red and purple, the \emph{right singular
vectors}) are the only ones that come out still perpendicular, landing on the
ellipse's two axes. Their output lengths $\sigma_1 \ge \sigma_2$ are the
\emph{singular values}.}
\label{ch29:fig-circle-to-ellipse}
\end{figure}

Here is the observation everything hinges on. As you sweep
an input unit vector around the circle, its image sweeps around the ellipse. Most
input directions get knocked off square -- two perpendicular inputs usually come
out \emph{not} perpendicular, because the matrix shears as well as stretches. But
there is always (at least) one special pair of perpendicular input directions
that stays perpendicular on the way out. Those special inputs land exactly on the
long axis and short axis of the ellipse.

\begin{intuitionbox}
The whole SVD lives in this one picture. A matrix turns the unit circle into an
ellipse. The \textbf{right singular vectors} $\vec{v}_1, \vec{v}_2$ are the
perpendicular \emph{input} directions that survive as perpendicular. The
\textbf{left singular vectors} $\vec{u}_1, \vec{u}_2$ are the perpendicular
\emph{output} directions they land on -- the ellipse's axes. The
\textbf{singular values} $\sigma_1 \ge \sigma_2 \ge 0$ are how long those output
axes are. The matrix does nothing more exotic than: line up with the
$\vec{v}$'s, stretch by the $\sigma$'s, point along the $\vec{u}$'s.
\end{intuitionbox}

Why must the image always be an ellipse, and why must such a perpendicular pair
always exist? That is not obvious, and we will prove it. But the geometry should
already feel inevitable: an ellipse is just a stretched, rotated circle, and a
matrix is just the machine that stretches and rotates. Let us turn the picture
into algebra.

\begin{recoverybox}
\textbf{Pause and think:} What does the output ellipse look like if $\mat{A}$ is
a $2\times 2$ projection matrix onto the diagonal line $y = x$?

A projection squashes everything perpendicular to the line down to zero, while
vectors along the line are untouched. The unit circle collapses into a flat
line segment running from $-\frac{1}{\sqrt2}\colvec{1\\1}$ to
$+\frac{1}{\sqrt2}\colvec{1\\1}$. Its semi-major axis length is $\sigma_1 = 1$,
and its semi-minor axis length is $\sigma_2 = 0$. A ``flat ellipse'' is still an
ellipse!
\end{recoverybox}

\section{From the Picture to the Factorization}
\label{ch29:to-factorization}

Suppose we have found these aligned orthogonal directions. In the plane: an orthonormal pair
$\vec{v}_1, \vec{v}_2$ in the input space, an orthonormal pair $\vec{u}_1,
\vec{u}_2$ in the output space, and nonnegative numbers $\sigma_1, \sigma_2$, all
linked by

\[
  \mat{A}\vec{v}_1 = \sigma_1 \vec{u}_1, \qquad
  \mat{A}\vec{v}_2 = \sigma_2 \vec{u}_2.
\]

Read that as: ``the input axis $\vec{v}_i$ gets sent to the output axis
$\vec{u}_i$, stretched by $\sigma_i$.'' Now stack these two equations as columns
of a single matrix equation. Put the $\vec{v}$'s as columns of a matrix
$\mat{V}$, the $\vec{u}$'s as columns of $\mat{U}$:

\[
  \mat{A}\underbrace{\begin{bmatrix} \vec{v}_1 & \vec{v}_2 \end{bmatrix}}_{\mat{V}}
  = \underbrace{\begin{bmatrix} \vec{u}_1 & \vec{u}_2 \end{bmatrix}}_{\mat{U}}
    \underbrace{\begin{bmatrix} \sigma_1 & 0 \\ 0 & \sigma_2 \end{bmatrix}}_{\mat{\Sigma}}.
\]

The right side is correct because $\mat{U}\mat{\Sigma}$ has columns $\sigma_1
\vec{u}_1$ and $\sigma_2 \vec{u}_2$ -- exactly what $\mat{A}\vec{v}_1$ and
$\mat{A}\vec{v}_2$ are supposed to equal. So $\mat{A}\mat{V} = \mat{U}\mat{\Sigma}$.
Because the columns of $\mat{V}$ are mutually perpendicular unit vectors
(orthonormal), $\mat{V}$ is an \emph{orthogonal matrix} -- a square matrix
representing a rigid rotation or reflection that preserves all lengths and
angles, satisfying $\mat{V}\inv = \mat{V}\T$. Multiply on the right by
$\mat{V}\T$ and you have it:

\[
  \boxed{\;\mat{A} = \mat{U}\,\mat{\Sigma}\,\mat{V}\T.\;}
\]

\begin{focusbox}
The SVD is just the equations $\mat{A}\vec{v}_i = \sigma_i \vec{u}_i$ collected
into matrices. Read $\mat{A} = \mat{U}\mat{\Sigma}\mat{V}\T$ \textbf{right to
left} as three motions applied to an input vector $\vec{x}$:
\begin{enumerate}
  \item $\mat{V}\T$ \textbf{rotates} the input space so the right singular vectors $\vec{v}_i$ line up with the
        coordinate axes;
  \item $\mat{\Sigma}$ \textbf{stretches} each coordinate axis by its singular
        value (and, if rectangular, changes the dimension);
  \item $\mat{U}$ \textbf{rotates} the stretched result into the output space.
\end{enumerate}
Rotate, stretch, rotate. Every matrix, no exceptions.
\end{focusbox}

\begin{figure}[H]
\centering
\begin{tikzpicture}[scale=0.78,>=Stealth]
  % stage 0: unit circle with v1,v2
  \begin{scope}[xshift=0cm]
    \draw[gray!50] (0,0) circle (1);
    \draw[->,vecred,thick] (0,0)--(0.7,0.7);
    \draw[->,structurepurple,thick] (0,0)--(-0.7,0.7);
    \node[below] at (0,-1.3) {$\vec{x}$ on circle};
    \node at (0,1.6) {$\vec{v}_1,\vec{v}_2$};
  \end{scope}
  \draw[->,thick] (1.6,0)--(2.6,0) node[midway,above]{$\mat{V}\T$};
  % stage 1: rotated so v's on axes
  \begin{scope}[xshift=4.2cm]
    \draw[gray!50] (0,0) circle (1);
    \draw[->,vecred,thick] (0,0)--(1,0);
    \draw[->,structurepurple,thick] (0,0)--(0,1);
    \node[below] at (0,-1.3) {axes aligned};
    \node at (0,1.6) {$\vec{e}_1,\vec{e}_2$};
  \end{scope}
  \draw[->,thick] (5.8,0)--(6.8,0) node[midway,above]{$\mat{\Sigma}$};
  % stage 2: stretched ellipse, axis-aligned
  \begin{scope}[xshift=9.4cm]
    \draw[targetgreen!70] (0,0) ellipse (1.7 and 0.8);
    \draw[->,vecred,thick] (0,0)--(1.7,0);
    \draw[->,structurepurple,thick] (0,0)--(0,0.8);
    \node[below] at (0,-1.3) {stretched};
    \node at (0,1.6) {$\sigma_1,\sigma_2$};
  \end{scope}
  \draw[->,thick] (11.4,0)--(12.4,0) node[midway,above]{$\mat{U}$};
  % stage 3: rotated ellipse
  \begin{scope}[xshift=15.4cm]
    \draw[targetgreen,thick,rotate around={25:(0,0)}] (0,0) ellipse (1.7 and 0.8);
    \draw[->,vecred,thick] (0,0)--({1.7*cos(25)},{1.7*sin(25)});
    \draw[->,structurepurple,thick] (0,0)--({-0.8*sin(25)},{0.8*cos(25)});
    \node[below] at (0,-1.3) {$\mat{A}\vec{x}$ on ellipse};
    \node at (0,1.6) {$\vec{u}_1,\vec{u}_2$};
  \end{scope}
\end{tikzpicture}
\caption{The SVD as rotate--stretch--rotate. Reading $\mat{A} =
\mat{U}\mat{\Sigma}\mat{V}\T$ from right to left: $\mat{V}\T$ rotates the right singular
vectors $\vec{v}_i$ onto the coordinate axes, $\mat{\Sigma}$ stretches along those axes by
$\sigma_1,\sigma_2$, and $\mat{U}$ rotates the result into final position.}
\label{ch29:fig-rotate-stretch-rotate}
\end{figure}

Compare this to diagonalization, $\mat{A} = \mat{X}\mat{\Lambda}\mat{X}\inv$. There
the same basis $\mat{X}$ appears on both sides, but it is generally \emph{not}
orthogonal, so $\mat{X}\inv \ne \mat{X}\T$, and the picture is a skew stretch
along oblique eigen-axes. The SVD pays a small price -- it allows \emph{two}
different orthonormal bases, an input one $\mat{V}$ and an output one $\mat{U}$ --
and in exchange it buys enormous freedom: the bases are always orthonormal, the
stretches $\sigma_i$ are always real and nonnegative, and the matrix need not even
be square.

\begin{strangbox}
Strang loves to write the SVD as $\mat{A}\vec{v}_i = \sigma_i \vec{u}_i$ and call
$\mat{V}$ and $\mat{U}$ ``the right and left singular vector bases.'' If you have
already taken his course, you have seen him insist that the SVD ``completes'' the
four-subspaces picture: $\mat{V}$ gives orthonormal bases for the row space and
null space, $\mat{U}$ for the column space and left null space, and $\mat{\Sigma}$
records exactly how the row space is stretched onto the column space. If you are
reading this \emph{before} Strang, hold onto the circle-to-ellipse picture; it is
the one mental image that makes his blackboard derivation feel obvious.
\end{strangbox}

\section{The Formal Statement}
\label{ch29:formal-statement}

Let us write down precisely what we are claiming, for a general $m \times n$
matrix.

\begin{formalbox}
\begin{theorem}[Singular Value Decomposition]
\label{ch29:thm-svd}
Let $\mat{A}$ be any real $m \times n$ matrix. Then there exist
\begin{itemize}
  \item an orthogonal $m \times m$ matrix $\mat{U}$ (columns $\vec{u}_1,\dots,\vec{u}_m$),
  \item an orthogonal $n \times n$ matrix $\mat{V}$ (columns $\vec{v}_1,\dots,\vec{v}_n$),
  \item an $m \times n$ ``diagonal'' matrix $\mat{\Sigma}$ whose only nonzero
        entries are $\sigma_1 \ge \sigma_2 \ge \cdots \ge \sigma_r > 0$ in the
        top-left, where $r = \rank(\mat{A})$,
\end{itemize}
such that
\[
  \mat{A} = \mat{U}\,\mat{\Sigma}\,\mat{V}\T.
\]
The numbers $\sigma_i$ are the \emph{singular values}; the columns of $\mat{V}$
are the \emph{right singular vectors}; the columns of $\mat{U}$ are the
\emph{left singular vectors}. The singular values are uniquely determined by
$\mat{A}$.
\end{theorem}
\end{formalbox}

A word on shapes, since rectangular matrices trip everyone up at first. If
$\mat{A}$ is $m \times n$, then $\mat{\Sigma}$ is also $m \times n$ -- the same
shape as $\mat{A}$ -- but it is as diagonal as a rectangle can be. For a
$3 \times 2$ matrix,

\[
  \mat{\Sigma} = \begin{bmatrix} \sigma_1 & 0 \\ 0 & \sigma_2 \\ 0 & 0 \end{bmatrix},
\]

and for a $2 \times 3$ matrix,

\[
  \mat{\Sigma} = \begin{bmatrix} \sigma_1 & 0 & 0 \\ 0 & \sigma_2 & 0 \end{bmatrix}.
\]

The singular values march down the main diagonal; everything off it is zero.

\begin{trapbox}
\textbf{Singular values are not eigenvalues.} They are nonnegative by
definition, even when the matrix has negative or complex eigenvalues, even when
the matrix is not square (and so has no eigenvalues at all). For the rotation by
$90^\circ$, $\mat{R} = \begin{bmatrix} 0 & -1 \\ 1 & 0\end{bmatrix}$, the
eigenvalues are $\pm i$ -- but both singular values are $1$, because a rotation
stretches nothing: it sends the unit circle to itself. Eigenvalues measure
``invariant directions and their scaling''; singular values measure ``how much
the matrix stretches space.'' Do not confuse them.
\end{trapbox}

\begin{recoverybox}
\textbf{Pause and think:} If a $2\times 2$ diagonal matrix is $\mat{A} = \begin{bmatrix} 3 & 0 \\ 0 & -5 \end{bmatrix}$, what are its singular values?

Can a singular value be negative? No: singular values measure physical stretch lengths, so by definition $\sigma_i \ge 0$, sorted as $\sigma_1 \ge \sigma_2 \ge 0$. Here $\sigma_1 = 5$ and $\sigma_2 = 3$. The minus sign is absorbed into the rotation/reflection matrices: $\mat{A}\colvec{0\\1} = -5\colvec{0\\1} = 5\colvec{0\\-1}$, so the input axis $\vec{v}_2 = \colvec{0\\1}$ is sent to the output axis $\vec{u}_2 = \colvec{0\\-1}$.
\end{recoverybox}

\subsection{The economy (thin) SVD}

When $m$ and $n$ are very different -- say a $10000 \times 50$ data matrix -- the
full $\mat{U}$ would be a gigantic $10000 \times 10000$ matrix, almost all of
whose columns multiply zeros in $\mat{\Sigma}$. Wasteful. The \textbf{economy
SVD} keeps only the columns that matter. If $r = \rank(\mat{A})$, write

\[
  \mat{A} = \mat{U}_r\,\mat{\Sigma}_r\,\mat{V}_r\T,
\]

where $\mat{U}_r$ is $m \times r$ (its columns the first $r$ left singular
vectors), $\mat{\Sigma}_r = \diag(\sigma_1,\dots,\sigma_r)$ is $r \times r$ and
genuinely square, and $\mat{V}_r$ is $n \times r$. Even leaner is the
\textbf{outer-product form}, which we will use constantly:

\[
  \boxed{\;\mat{A} = \sum_{i=1}^{r} \sigma_i\, \vec{u}_i \vec{v}_i\T.\;}
\]

Each term $\sigma_i \vec{u}_i \vec{v}_i\T$ is a rank-one matrix -- an output
direction $\vec{u}_i$ ($m \times 1$) times an input direction $\vec{v}_i\T$ ($1 \times n$),
scaled by $\sigma_i$. The matrix is literally a \emph{sum of rank-one pieces}, ordered from
strongest to weakest (Figure~\ref{ch29:fig-outer-product}). This representation forms the foundation of low-rank matrix approximation and data compression.

\begin{figure}[H]
\centering
\begin{tikzpicture}[scale=0.88,>=Stealth]
  % A matrix
  \draw[thick,fill=gray!10] (0,0) rectangle (2.0,2.6);
  \node at (1.0,1.3) {$\mat{A}$};
  \node[below,font=\small] at (1.0,-0.1) {$m \times n$};

  \node at (2.5,1.3) {$=$};

  % Term 1: sigma1 * u1 * v1^T
  \node at (3.1,1.3) {$\sigma_1$};
  % u1 (m x 1 column)
  \draw[thick,targetgreen,fill=targetgreen!15] (3.5,0) rectangle (3.9,2.6);
  \node[targetgreen,font=\small] at (3.7,1.3) {$\vec{u}_1$};
  \node[font=\footnotesize] at (4.2,1.3) {$\times$};
  % v1^T (1 x n row)
  \draw[thick,sourceblue,fill=sourceblue!15] (4.5,2.1) rectangle (6.5,2.65);
  \node[sourceblue,font=\small] at (5.5,2.375) {$\vec{v}_1\T$};
  \node[below,font=\footnotesize,align=center] at (5.0,-0.1) {rank $1$\\(weight $\sigma_1$)};

  \node at (7.0,1.3) {$+$};

  % Term 2: sigma2 * u2 * v2^T
  \node at (7.6,1.3) {$\sigma_2$};
  % u2
  \draw[thick,targetgreen,fill=targetgreen!15] (8.0,0) rectangle (8.4,2.6);
  \node[targetgreen,font=\small] at (8.2,1.3) {$\vec{u}_2$};
  \node[font=\footnotesize] at (8.7,1.3) {$\times$};
  % v2^T
  \draw[thick,sourceblue,fill=sourceblue!15] (9.0,2.1) rectangle (11.0,2.65);
  \node[sourceblue,font=\small] at (10.0,2.375) {$\vec{v}_2\T$};
  \node[below,font=\footnotesize,align=center] at (9.5,-0.1) {rank $1$\\(weight $\sigma_2$)};

  \node at (11.6,1.3) {$+ \;\cdots\; +$};

  % Term r: sigma_r * ur * vr^T
  \node at (12.7,1.3) {$\sigma_r$};
  \draw[thick,targetgreen,fill=targetgreen!15] (13.1,0) rectangle (13.5,2.6);
  \node[targetgreen,font=\small] at (13.3,1.3) {$\vec{u}_r$};
  \node[font=\footnotesize] at (13.8,1.3) {$\times$};
  \draw[thick,sourceblue,fill=sourceblue!15] (14.1,2.1) rectangle (16.1,2.65);
  \node[sourceblue,font=\small] at (15.1,2.375) {$\vec{v}_r\T$};
  \node[below,font=\footnotesize,align=center] at (14.6,-0.1) {rank $1$\\(weight $\sigma_r$)};
\end{tikzpicture}
\caption{The outer-product expansion decomposes $\mat{A}$ into a sum of $r$ rank-one matrices. Each slice is an outer product $\vec{u}_i \vec{v}_i\T$ (an $m\times 1$ output column times a $1\times n$ input row) scaled by the singular value $\sigma_i$. Truncating after $k$ slices discards only the weakest remaining patterns.}
\label{ch29:fig-outer-product}
\end{figure}

\section{Where the Pieces Come From: $\mat{A}\T\mat{A}$ and $\mat{A}\mat{A}\T$}
\label{ch29:ata-aat}

So far we have assumed these aligned orthonormal directions exist. Now we construct them explicitly
and prove the existence of the SVD. The key step is to form symmetric matrices by multiplying
$\mat{A}$ with its transpose: both $\mat{A}\T\mat{A}$ and $\mat{A}\mat{A}\T$ are symmetric and
positive semidefinite, so the spectral theorem applies directly to them.

Suppose for the moment that $\mat{A} = \mat{U}\mat{\Sigma}\mat{V}\T$ really holds.
Then watch what $\mat{A}\T\mat{A}$ becomes, using $\mat{U}\T\mat{U} = \mat{I}$:

\[
  \mat{A}\T\mat{A}
  = (\mat{U}\mat{\Sigma}\mat{V}\T)\T(\mat{U}\mat{\Sigma}\mat{V}\T)
  = \mat{V}\mat{\Sigma}\T\mat{U}\T\mat{U}\mat{\Sigma}\mat{V}\T
  = \mat{V}\,(\mat{\Sigma}\T\mat{\Sigma})\,\mat{V}\T.
\]

But $\mat{\Sigma}\T\mat{\Sigma}$ is the diagonal matrix
$\diag(\sigma_1^2,\dots,\sigma_n^2)$. So this last line \emph{is} an
orthogonal diagonalization of $\mat{A}\T\mat{A}$:

\begin{focusbox}
The right singular vectors $\vec{v}_i$ are the \textbf{eigenvectors of
$\mat{A}\T\mat{A}$}, and the squared singular values $\sigma_i^2$ are its
\textbf{eigenvalues}:
\[
  \mat{A}\T\mat{A}\,\vec{v}_i = \sigma_i^2\, \vec{v}_i.
\]
Symmetrically, the left singular vectors $\vec{u}_i$ are the eigenvectors of
$\mat{A}\mat{A}\T$ with the \emph{same} eigenvalues $\sigma_i^2$:
\[
  \mat{A}\mat{A}\T\,\vec{u}_i = \sigma_i^2\, \vec{u}_i.
\]
\end{focusbox}

This structural relationship provides both a constructive proof and a practical recipe. Let us see why
it must work, then turn it into a procedure.

\subsection{Why the recipe is forced (a real proof)}

We build the SVD from scratch, with no prior assumption that it exists. Start with
$\mat{A}\T\mat{A}$. It is $n\times n$, symmetric, and \emph{positive
semidefinite}: for any $\vec{x}$,
\[
  \vec{x}\T(\mat{A}\T\mat{A})\vec{x} = (\mat{A}\vec{x})\T(\mat{A}\vec{x})
  = \norm{\mat{A}\vec{x}}^2 \ge 0.
\]
By the spectral theorem (the symmetric chapter), $\mat{A}\T\mat{A}$ has an
orthonormal basis of eigenvectors $\vec{v}_1,\dots,\vec{v}_n$ with real
eigenvalues $\lambda_1 \ge \cdots \ge \lambda_n$, and because it is positive
semidefinite, every $\lambda_i \ge 0$. Define the singular values as the
square roots,
\[
  \sigma_i := \sqrt{\lambda_i} \ge 0,
\]
and say there are $r$ positive ones: $\sigma_1 \ge \cdots \ge \sigma_r > 0$ and
$\sigma_{r+1} = \cdots = 0$.

Now \emph{define} the left singular vectors, for $i = 1,\dots,r$, by normalizing
the images $\mat{A}\vec{v}_i$:
\[
  \vec{u}_i := \frac{1}{\sigma_i}\,\mat{A}\vec{v}_i.
\]
Two things to check: that these $\vec{u}_i$ are orthonormal, and that the leftover
directions ($i > r$) really do contribute nothing.

\emph{Orthonormality.} Compute the dot product of two of them:
\[
  \vec{u}_i\T\vec{u}_j
  = \frac{1}{\sigma_i\sigma_j}(\mat{A}\vec{v}_i)\T(\mat{A}\vec{v}_j)
  = \frac{1}{\sigma_i\sigma_j}\,\vec{v}_i\T(\mat{A}\T\mat{A})\vec{v}_j
  = \frac{1}{\sigma_i\sigma_j}\,\vec{v}_i\T(\sigma_j^2\,\vec{v}_j)
  = \frac{\sigma_j}{\sigma_i}\,\vec{v}_i\T\vec{v}_j.
\]
Since the $\vec{v}$'s are orthonormal, $\vec{v}_i\T\vec{v}_j$ is $1$ when $i=j$
and $0$ otherwise. So $\vec{u}_i\T\vec{u}_j = 1$ when $i=j$ (the $\sigma$'s
cancel) and $0$ otherwise. The $\vec{u}_i$ are orthonormal.

\emph{The leftovers vanish.} For $i > r$ we have $\sigma_i = 0$, meaning
$\lambda_i = 0$, so
\[
  \norm{\mat{A}\vec{v}_i}^2 = \vec{v}_i\T(\mat{A}\T\mat{A})\vec{v}_i
  = \lambda_i \norm{\vec{v}_i}^2 = 0,
\]
hence $\mat{A}\vec{v}_i = \vec{0}$. These directions are crushed to nothing; they
are the null space of $\mat{A}$.

Finally, extend $\vec{u}_1,\dots,\vec{u}_r$ to a full orthonormal basis
$\vec{u}_1,\dots,\vec{u}_m$ of $\R^m$ (always possible -- just complete the
basis). By construction $\mat{A}\vec{v}_i = \sigma_i\vec{u}_i$ holds for
\emph{every} $i$ (for $i \le r$ by definition of $\vec{u}_i$; for $i > r$ both
sides are zero). Collect into matrices: $\mat{A}\mat{V} = \mat{U}\mat{\Sigma}$, so
$\mat{A} = \mat{U}\mat{\Sigma}\mat{V}\T$. \qed

\begin{aestheticbox}
Look at what just happened. We never assumed $\mat{A}$ was square, never assumed
it was invertible, never touched a complex number. Because $\mat{A}\T\mat{A}$ is
always symmetric and positive semidefinite, applying the spectral theorem to
$\mat{A}\T\mat{A}$ yields an orthogonal decomposition for \emph{any} matrix
$\mat{A}$.
\end{aestheticbox}

\begin{trapbox}
Use $\mat{A}\T\mat{A}$ to get the $\vec{v}$'s, then get the $\vec{u}$'s
\emph{from} $\vec{u}_i = \mat{A}\vec{v}_i / \sigma_i$ -- do \emph{not} find the
$\vec{u}$'s independently as eigenvectors of $\mat{A}\mat{A}\T$ and pair them up
arbitrarily. Eigenvectors are only defined up to sign (and up to rotation within
a repeated eigenspace). If you compute $\mat{U}$ and $\mat{V}$ separately, you
can easily get signs that disagree, so that $\mat{U}\mat{\Sigma}\mat{V}\T$ comes
out as $-\mat{A}$ in some columns. The link $\mat{A}\vec{v}_i = \sigma_i\vec{u}_i$
fixes the signs consistently.
\end{trapbox}

\section{A Fully Worked $2 \times 2$ Example}
\label{ch29:worked-2x2}

Let us compute a $2\times 2$ SVD from start to finish. Take

\[
  \mat{A} = \begin{bmatrix} 3 & 0 \\ 4 & 5 \end{bmatrix}.
\]

\textbf{Step 1: Form $\mat{A}\T\mat{A}$.}
\[
  \mat{A}\T\mat{A}
  = \begin{bmatrix} 3 & 4 \\ 0 & 5 \end{bmatrix}
    \begin{bmatrix} 3 & 0 \\ 4 & 5 \end{bmatrix}
  = \begin{bmatrix} 25 & 20 \\ 20 & 25 \end{bmatrix}.
\]

\textbf{Step 2: Eigenvalues of $\mat{A}\T\mat{A}$.} For the symmetric matrix
$\begin{bmatrix} 25 & 20 \\ 20 & 25\end{bmatrix}$, the eigenvalues of
$\begin{bmatrix} a & b \\ b & a\end{bmatrix}$ are $a+b$ and $a-b$ (the vectors
$\colvec{1\\1}$ and $\colvec{1\\-1}$ are always eigenvectors of such a matrix).
So
\[
  \lambda_1 = 25 + 20 = 45, \qquad \lambda_2 = 25 - 20 = 5.
\]
The singular values are the square roots:
\[
  \sigma_1 = \sqrt{45} = 3\sqrt{5} \approx 6.708, \qquad
  \sigma_2 = \sqrt{5} \approx 2.236.
\]

\textbf{Step 3: Right singular vectors (the $\vec{v}$'s).} The eigenvectors are
$\colvec{1\\1}$ and $\colvec{1\\-1}$; normalize them:
\[
  \vec{v}_1 = \frac{1}{\sqrt2}\colvec{1\\1}, \qquad
  \vec{v}_2 = \frac{1}{\sqrt2}\colvec{1\\-1}.
\]

\textbf{Step 4: Left singular vectors (the $\vec{u}$'s)} via $\vec{u}_i =
\mat{A}\vec{v}_i / \sigma_i$. First
\[
  \mat{A}\vec{v}_1 = \frac{1}{\sqrt2}\begin{bmatrix} 3 & 0 \\ 4 & 5 \end{bmatrix}\colvec{1\\1}
  = \frac{1}{\sqrt2}\colvec{3\\9},
  \quad\Rightarrow\quad
  \vec{u}_1 = \frac{1}{\sigma_1}\cdot\frac{1}{\sqrt2}\colvec{3\\9}
  = \frac{1}{3\sqrt5}\cdot\frac{1}{\sqrt2}\colvec{3\\9}
  = \frac{1}{\sqrt{10}}\colvec{1\\3}.
\]
Check its length: $\frac{1}{\sqrt{10}}\sqrt{1+9} = 1$. Good. Then
\[
  \mat{A}\vec{v}_2 = \frac{1}{\sqrt2}\begin{bmatrix} 3 & 0 \\ 4 & 5 \end{bmatrix}\colvec{1\\-1}
  = \frac{1}{\sqrt2}\colvec{3\\-1},
  \quad\Rightarrow\quad
  \vec{u}_2 = \frac{1}{\sqrt5}\cdot\frac{1}{\sqrt2}\colvec{3\\-1}
  = \frac{1}{\sqrt{10}}\colvec{3\\-1}.
\]
Notice $\vec{u}_1\T\vec{u}_2 = \frac{1}{10}(1\cdot 3 + 3\cdot(-1)) = 0$: the
output axes really are perpendicular, exactly as the picture promised.

\textbf{Step 5: Assemble.}
\[
  \mat{A}
  = \underbrace{\frac{1}{\sqrt{10}}\begin{bmatrix} 1 & 3 \\ 3 & -1 \end{bmatrix}}_{\mat{U}}
    \underbrace{\begin{bmatrix} 3\sqrt5 & 0 \\ 0 & \sqrt5 \end{bmatrix}}_{\mat{\Sigma}}
    \underbrace{\frac{1}{\sqrt2}\begin{bmatrix} 1 & 1 \\ 1 & -1 \end{bmatrix}}_{\mat{V}\T}.
\]

\textbf{Step 6: Sanity check} by multiplying $\mat{U}\mat{\Sigma}$ first:
\[
  \mat{U}\mat{\Sigma}
  = \frac{1}{\sqrt{10}}\begin{bmatrix} 1 & 3 \\ 3 & -1 \end{bmatrix}
    \begin{bmatrix} 3\sqrt5 & 0 \\ 0 & \sqrt5 \end{bmatrix}
  = \frac{1}{\sqrt{10}}\begin{bmatrix} 3\sqrt5 & 3\sqrt5 \\ 9\sqrt5 & -\sqrt5 \end{bmatrix}.
\]
Then times $\mat{V}\T$, using $\frac{\sqrt5}{\sqrt{10}}\cdot\frac{1}{\sqrt2} =
\frac{\sqrt5}{\sqrt{20}} = \frac{\sqrt5}{2\sqrt5} = \frac12$:
\[
  \mat{U}\mat{\Sigma}\mat{V}\T
  = \frac12\begin{bmatrix} 3 & 3 \\ 9 & -1 \end{bmatrix}
           \begin{bmatrix} 1 & 1 \\ 1 & -1 \end{bmatrix}
  = \frac12\begin{bmatrix} 6 & 0 \\ 8 & 10 \end{bmatrix}
  = \begin{bmatrix} 3 & 0 \\ 4 & 5 \end{bmatrix} = \mat{A}. \checkmark
\]

It closes perfectly. Notice also that $\sigma_1\sigma_2 = 3\sqrt5\cdot\sqrt5 = 15
= \det\mat{A}$ (since $\det\mat{A} = 3\cdot5 - 0\cdot4 = 15$). That is no
accident -- see the next box.

\begin{labbox}
Compute the SVD of $\mat{A} = \begin{bmatrix} 4 & 0 \\ 3 & -5\end{bmatrix}$ by
hand, following the six steps above. You should find $\mat{A}\T\mat{A} =
\begin{bmatrix} 25 & -15 \\ -15 & 25\end{bmatrix}$, eigenvalues $40$ and $10$, and
singular values $\sigma_1 = \sqrt{40} = 2\sqrt{10}$, $\sigma_2 = \sqrt{10}$. Find
the $\vec{v}$'s, push them through $\mat{A}$, and check that the two $\vec{u}$'s
come out orthogonal. Then verify $\sigma_1\sigma_2 = \abs{\det\mat{A}} = 20$.
\end{labbox}

\begin{connectionbox}
\textbf{Singular values, determinant, and volume.} For a square $\mat{A}$,
\[
  \abs{\det\mat{A}} = \sigma_1\sigma_2\cdots\sigma_n.
\]
Here is the slick reason. Determinants multiply, and orthogonal matrices have
determinant $\pm 1$ (they preserve volume), so
$\abs{\det\mat{A}} = \abs{\det\mat{U}}\,\abs{\det\mat{\Sigma}}\,\abs{\det\mat{V}\T}
= 1\cdot(\sigma_1\cdots\sigma_n)\cdot 1$. Geometrically: the unit cube has
volume $1$; rotating ($\mat{U}$, $\mat{V}\T$) does not change volume; only the
stretch $\mat{\Sigma}$ does, multiplying volume by $\sigma_1\cdots\sigma_n$. The
absolute value of the determinant is exactly the volume of the output box -- and
the SVD shows you which directions did the stretching.
\end{connectionbox}

\section{The SVD and the Four Fundamental Subspaces}
\label{ch29:four-subspaces}

Recall the four fundamental subspaces of an $m\times n$ matrix $\mat{A}$ of rank $r$:
\begin{itemize}
  \item the \textbf{column space} $\col(\mat{A}) \subseteq \R^m$, dimension $r$;
  \item the \textbf{row space} $\row(\mat{A}) = \col(\mat{A}\T) \subseteq \R^n$, dimension $r$;
  \item the \textbf{null space} $\nullsp(\mat{A}) \subseteq \R^n$, dimension $n-r$;
  \item the \textbf{left null space} $\nullsp(\mat{A}\T) \subseteq \R^m$, dimension $m-r$.
\end{itemize}
The SVD hands you a clean orthonormal basis for all four at once. Split the
singular vectors at the rank $r$:

\begin{focusbox}
\begin{itemize}
  \item $\vec{v}_1,\dots,\vec{v}_r$ (positive singular values) are an orthonormal
        basis for the \textbf{row space}.
  \item $\vec{v}_{r+1},\dots,\vec{v}_n$ (zero singular values) are an orthonormal
        basis for the \textbf{null space}.
  \item $\vec{u}_1,\dots,\vec{u}_r$ are an orthonormal basis for the
        \textbf{column space}.
  \item $\vec{u}_{r+1},\dots,\vec{u}_m$ are an orthonormal basis for the
        \textbf{left null space}.
\end{itemize}
And $\mat{A}$ maps each $\vec{v}_i$ ($i\le r$) onto $\sigma_i \vec{u}_i$: the row
space gets rotated and stretched cleanly onto the column space.
\end{focusbox}

Why? We already saw that $\vec{v}_i$ with $\sigma_i = 0$ satisfy $\mat{A}\vec{v}_i
= \vec{0}$, so they live in the null space; there are $n - r$ of them, which is
exactly the null space's dimension, so they are a basis for it. The remaining
$\vec{v}_1,\dots,\vec{v}_r$ are orthogonal to the null space, and the orthogonal
complement of the null space is the row space (fundamental theorem of linear
algebra) -- so they are an orthonormal basis for the row space. For the
$\vec{u}$'s: each $\vec{u}_i = \mat{A}\vec{v}_i/\sigma_i$ ($i\le r$) is a column of
$\mat{A}$ times a vector, hence in $\col(\mat{A})$; there are $r$ of them and they
are orthonormal, so they are a basis for the $r$-dimensional column space. The
leftover $\vec{u}_{r+1},\dots,\vec{u}_m$ are orthogonal to the column space, i.e.
in the left null space.

\begin{figure}[H]
\centering
\begin{tikzpicture}[scale=0.9,>=Stealth]
  % input space R^n
  \node at (-2.5,3.6) {\textbf{input space} $\R^n$};
  \draw[sourceblue,thick,fill=sourceblue!8] (-4.4,1.6) rectangle (-0.6,3.0);
  \node[align=center] at (-2.5,2.3) {row space\\$\spn(\vec{v}_1,\dots,\vec{v}_r)$};
  \draw[trapred,thick,fill=trapred!8] (-4.4,-0.2) rectangle (-0.6,1.2);
  \node[align=center] at (-2.5,0.5) {null space\\$\spn(\vec{v}_{r+1},\dots,\vec{v}_n)$};
  % output space R^m
  \node at (5.0,3.6) {\textbf{output space} $\R^m$};
  \draw[targetgreen,thick,fill=targetgreen!8] (3.1,1.6) rectangle (6.9,3.0);
  \node[align=center] at (5.0,2.3) {column space\\$\spn(\vec{u}_1,\dots,\vec{u}_r)$};
  \draw[quotientorange,thick,fill=quotientorange!8] (3.1,-0.2) rectangle (6.9,1.2);
  \node[align=center] at (5.0,0.5) {left null space\\$\spn(\vec{u}_{r+1},\dots,\vec{u}_m)$};
  % origin in output space
  \fill (3.1,1.4) circle (2.5pt) node[right=3pt,font=\small]{$\vec{0}\in\R^m$};
  % forward action on row space
  \draw[->,thick] (-0.5,2.4) -- (3.0,2.4) node[midway,above]{$\mat{A}\vec{v}_i = \sigma_i\vec{u}_i$};
  % reverse action of pseudoinverse
  \draw[->,thick,dashed,sourceblue] (3.0,1.9) -- (-0.5,1.9)
        node[midway,below,font=\small]{$\mat{A}^+\vec{u}_i = \tfrac{1}{\sigma_i}\vec{v}_i$};
  % action on null space
  \draw[->,thick,trapred] (-0.5,0.5) .. controls (1.2,0.3) .. (3.0,1.38)
        node[midway,below=2pt,font=\small]{$\mat{A}\vec{v} = \vec{0}$};
\end{tikzpicture}
\caption{The SVD completes the four-subspaces picture. The right singular vectors
split the input space $\R^n$ into row space (top, stretched onto the column space
by $\sigma_i$) and null space (bottom, crushed to $\vec{0}\in\R^m$). The left singular vectors
split the output space $\R^m$ into column space (the reachable image) and left null space
(perpendicular to the image, never reached). The pseudoinverse $\mat{A}^+$ runs the top
stretch in reverse, dividing each output axis by $\sigma_i$.}
\label{ch29:fig-four-subspaces}
\end{figure}

\begin{recoverybox}
\textbf{Pause and think:} In Figure~\ref{ch29:fig-four-subspaces}, what happens to the null space and left null space if $\mat{A}$ is an invertible $n\times n$ matrix?

If $\mat{A}$ is invertible, its rank is $r = n = m$. The null space and left null space have dimension $n - r = 0$: they both shrink to just the origin $\{\vec{0}\}$. The row space fills all of $\R^n$, the column space fills all of $\R^m$, and every singular value is strictly positive ($\sigma_n > 0$). The transformation is an invertible, one-to-one stretch (up to rotation) along all $n$ axes.
\end{recoverybox}

\begin{strangbox}
This is the figure Strang calls ``the big picture of linear algebra,'' now drawn
with \emph{orthonormal} axes. In his earlier chapters the four subspaces are
there, but the bases are whatever elimination happened to give you -- not
orthogonal, not unit length. The SVD is the moment those four rooms get
perpendicular coordinate axes and the action of $\mat{A}$ becomes a simple
diagonal stretch from one to the other. Many readers say the four subspaces
finally ``clicked'' only when they saw the SVD. If that is you, you are in good
company.
\end{strangbox}

\section{Best Low-Rank Approximation: Eckart--Young}
\label{ch29:eckart-young}

Now the payoff that makes the SVD indispensable in the data age. Recall the
outer-product form,
\[
  \mat{A} = \sigma_1 \vec{u}_1\vec{v}_1\T + \sigma_2 \vec{u}_2\vec{v}_2\T
            + \cdots + \sigma_r \vec{u}_r\vec{v}_r\T.
\]
The terms are sorted: $\sigma_1 \ge \sigma_2 \ge \cdots$, so the \emph{first}
term carries the most ``energy,'' the second the next most, and so on. A natural
idea: keep only the top $k$ terms and throw away the rest. Define
\[
  \mat{A}_k := \sum_{i=1}^{k} \sigma_i\,\vec{u}_i\vec{v}_i\T.
\]
This $\mat{A}_k$ has rank $k$ (a sum of $k$ rank-one pieces). It is an
approximation to $\mat{A}$ built from the $k$ strongest directions.

To say that $\mat{A}_k$ is the ``best'' approximation, we need a way to
measure the error matrix $\mat{A} - \mat{B}$. Two standard matrix norms answer
this:
\begin{itemize}
  \item The \textbf{spectral norm} (or operator 2-norm) $\norm{\mat{M}}_2 := \max_{\norm{\vec{x}}=1}\norm{\mat{M}\vec{x}}$
        measures the \emph{worst-case stretch} -- the largest factor by which $\mat{M}$
        can lengthen any unit input vector.
  \item The \textbf{Frobenius norm} $\norm{\mat{M}}_F := \sqrt{\sum_{i,j} M_{ij}^2}$
        measures the \emph{total mass} -- the square root of the sum of squared
        entries, treating the entire matrix like a single long Euclidean vector in $\R^{mn}$.
\end{itemize}
The theorem of Eckart and Young says: \emph{this is the best rank-$k$
approximation there is}, under both measures simultaneously. Nothing of rank $k$ comes closer.

\begin{formalbox}
\begin{theorem}[Eckart--Young, 1936]
\label{ch29:thm-eckart-young}
Let $\mat{A} = \sum_{i=1}^r \sigma_i \vec{u}_i\vec{v}_i\T$ and let $\mat{A}_k$ be
its truncation to the top $k$ terms. Then for \emph{every} matrix $\mat{B}$ of
rank at most $k$,
\[
  \norm{\mat{A} - \mat{A}_k} \le \norm{\mat{A} - \mat{B}},
\]
in both the spectral norm and the Frobenius norm. The error left behind is
exactly
\[
  \norm{\mat{A} - \mat{A}_k}_{2} = \sigma_{k+1}, \qquad
  \norm{\mat{A} - \mat{A}_k}_{F} = \sqrt{\sigma_{k+1}^2 + \cdots + \sigma_r^2}.
\]
\end{theorem}
\end{formalbox}

This error formula explains why the SVD is so widely used in data
analysis and scientific computing. The error of the best rank-$k$ approximation is governed by the
\emph{discarded singular values}. If $\sigma_{k+1}, \sigma_{k+2}, \dots$ are tiny,
then $\mat{A}_k$ is nearly equal to $\mat{A}$ even though it is enormously
simpler. A $1000\times 1000$ matrix whose singular values plummet after, say,
$\sigma_{20}$ can be captured almost perfectly by $20$ rank-one pieces -- that is
$20\cdot(1000 + 1000)$ numbers instead of $1{,}000{,}000$. Compression, denoising,
and dimensionality reduction are all this one move.
Figure~\ref{ch29:fig-sv-spectrum} shows the typical shape that makes it work: a
few dominant singular values and a long, negligible tail.

\begin{figure}[H]
\centering
\begin{tikzpicture}[scale=0.9,>=Stealth]
  \begin{axis}[
      width=10.5cm, height=5.2cm,
      xlabel={index $i$}, ylabel={$\sigma_i$},
      xmin=0.3, xmax=12.7, ymin=0, ymax=10.5,
      xtick={1,2,...,12}, ytick={0,2,4,6,8,10},
      ybar, bar width=8pt,
      axis lines=left, enlarge x limits=0.04,
      tick label style={font=\small}, label style={font=\small},
    ]
    \addplot[draw=sourceblue,fill=sourceblue!55] coordinates {
      (1,9.8) (2,7.1) (3,4.4) (4,2.7)
      (5,0.55) (6,0.42) (7,0.34) (8,0.28)
      (9,0.22) (10,0.18) (11,0.13) (12,0.09) };
    \draw[trapred,thick,dashed] (axis cs:4.5,0) -- (axis cs:4.5,10.2);
    \node[trapred,anchor=west,font=\small] at (axis cs:4.6,9.3)
      {keep $k=4$ $\;\to\;$ tail is the error};
  \end{axis}
\end{tikzpicture}
\caption{A typical singular-value spectrum: a few large values, then a long
low tail. Truncating at $k=4$ keeps the four big rank-one pieces; by
Eckart--Young the leftover error is governed entirely by the discarded tail
$\sigma_5,\sigma_6,\dots$, which here is tiny. This sharp drop-off is why real
data matrices compress so well.}
\label{ch29:fig-sv-spectrum}
\end{figure}

\begin{intuitionbox}
Think of each rank-one term $\sigma_i\vec{u}_i\vec{v}_i\T$ as an elementary
pattern weighted by $\sigma_i$. The SVD writes $\mat{A}$ as a sum of mutually
orthogonal patterns ordered from largest weight to smallest, and truncating at
$k$ keeps the $k$ strongest patterns while discarding the smallest ones. Because the directions are
orthogonal, dropping the small ones costs you \emph{only} their own energy --
nothing leaks. That orthogonality is why truncating at rank $k$ is provably
optimal, minimizing total reconstruction error.
\end{intuitionbox}

\begin{recoverybox}
\textbf{Pause and think:} Suppose an $m \times n$ matrix $\mat{A}$ already has rank 4. What is the error of its rank-4 SVD truncation $\mat{A}_4$?

Since the rank is 4, all singular values past the fourth are zero: $\sigma_5 = \sigma_6 = \cdots = 0$. By Eckart--Young, the spectral error is $\sigma_5 = 0$ and the Frobenius error is $\sqrt{\sigma_5^2 + \cdots} = 0$. The approximation is completely exact: $\mat{A}_4 = \mat{A}$. Truncation only discards information when $k < r$.
\end{recoverybox}

\subsection{Why the Frobenius error is what it is}

We will not prove full optimality over all $\mat{B}$ (that takes a careful
argument), but we can see \emph{why the truncation error is the tail of the
singular values}, which is the part you use in practice. The Frobenius norm
squared is the sum of squares of all entries, and a beautiful fact (provable by
expanding $\norm{\mat{M}}_F^2 = \tr(\mat{M}\T\mat{M})$ and using orthonormality)
is that it equals the sum of squared singular values:
\[
  \norm{\mat{A}}_F^2 = \sigma_1^2 + \sigma_2^2 + \cdots + \sigma_r^2.
\]
The discarded part $\mat{A} - \mat{A}_k = \sum_{i=k+1}^r \sigma_i\vec{u}_i\vec{v}_i\T$
is itself a (sub-)SVD with singular values $\sigma_{k+1},\dots,\sigma_r$, so by
the same fact
\[
  \norm{\mat{A} - \mat{A}_k}_F^2 = \sigma_{k+1}^2 + \cdots + \sigma_r^2.
\]
The orthogonality of the $\vec{u}$'s and $\vec{v}$'s is what makes the cross terms
vanish and the energies simply add. \emph{No information overlaps between the
pieces}, so you can account for them one at a time.

\begin{applicationbox}
\textbf{Image compression and the data that runs machine learning.} Store a
grayscale photo as an $m\times n$ matrix of pixel brightnesses and take its SVD.
Natural images are typically low-rank: a few dozen singular values capture nearly
all of what your eye sees, and the rest are essentially texture noise. Keeping the
top $k$ terms, $\mat{A}_k = \sum_{i=1}^k \sigma_i\vec{u}_i\vec{v}_i\T$, gives a
recognizable image stored in $k(m+n+1)$ numbers instead of $mn$ -- often a $10\times$
or $20\times$ saving with no visible loss.

The same move is everywhere in ML. \textbf{Latent semantic analysis} takes a
huge term-by-document matrix and truncates its SVD to discover a few dozen
``topics.'' \textbf{Recommender systems} treat the (mostly empty) user-by-item
ratings matrix as approximately low rank and use a truncated SVD to fill in the
blanks -- the famous Netflix-prize approach. \textbf{Modern large language
models} are compressed and fine-tuned with \emph{low-rank adapters} (LoRA), which
add a trainable rank-$k$ update $\mat{B}\mat{C}$ -- exactly the
$\sum \sigma_i \vec{u}_i\vec{v}_i\T$ shape -- to giant weight matrices, cutting the
number of trainable parameters by orders of magnitude. In each setting, the
Eckart--Young theorem guarantees that truncating the SVD gives the closest
rank-$k$ matrix in both the spectral and Frobenius norms.
\end{applicationbox}

\begin{walkbox}
Take a walk and turn this over: a rank-one matrix $\sigma\vec{u}\vec{v}\T$ is the
simplest non-trivial matrix there is -- one output pattern $\vec{u}$, one input
pattern $\vec{v}$, one gain $\sigma$. The SVD says \emph{every} matrix is a stack
of these, ranked by importance and mutually perpendicular. Why should rank-one
pieces, the humblest building blocks, be enough to reconstruct any matrix
whatsoever? And why is ``perpendicular and ranked'' the arrangement that makes
truncation optimal rather than merely convenient?
\end{walkbox}

\section{The Pseudoinverse: Solving the Unsolvable}
\label{ch29:pseudoinverse}

What does it mean to ``invert'' a matrix that has no inverse -- a rectangular one,
or a singular one? The SVD gives the right answer. To invert
$\mat{A} = \mat{U}\mat{\Sigma}\mat{V}\T$ you would like to undo each step:
$\mat{U}$ is orthogonal so its inverse is $\mat{U}\T$; same for $\mat{V}$. The only
trouble is $\mat{\Sigma}$, which you cannot invert if some singular values are
zero (you cannot divide by zero) or if it is not square. So you do the only
sensible thing: invert what you can and leave the rest at zero. Define
$\mat{\Sigma}^+$ by replacing each nonzero $\sigma_i$ with $1/\sigma_i$ and
transposing the shape. Then the \textbf{Moore--Penrose pseudoinverse} is
\[
  \boxed{\;\mat{A}^+ = \mat{V}\,\mat{\Sigma}^+\,\mat{U}\T.\;}
\]

\begin{focusbox}
The pseudoinverse $\mat{A}^+$ reverses $\mat{A}$ on the directions it can
(the row/column spaces, dividing by $\sigma_i$) and sends the unreachable
directions to zero. When $\mat{A}$ is genuinely invertible, $\mat{A}^+ =
\mat{A}\inv$ exactly. When it is not, $\mat{A}^+\vec{b}$ delivers the
\emph{least-squares, minimum-norm} solution of $\mat{A}\vec{x} = \vec{b}$ -- the
best you could possibly hope for.
\end{focusbox}

\begin{figure}[H]
\centering
\begin{tikzpicture}[scale=0.9,>=Stealth]
  % input plane R^2
  \begin{scope}[xshift=0cm]
    \draw[->,gray!55] (-1.7,0)--(1.9,0) node[right]{};
    \draw[->,gray!55] (0,-1.7)--(0,1.9) node[above]{};
    \draw[sourceblue,thick] (0,0) circle (1);
    \draw[->,vecred,very thick] (0,0)--(1,0) node[below right,vecred]{$\vec{v}_1$};
    \draw[->,structurepurple,very thick] (0,0)--(0,1) node[above left,structurepurple]{$\vec{v}_2$};
    \node at (0,-2.2) {input $\R^2$};
  \end{scope}
  \draw[->,thick] (2.2,0)--(3.6,0) node[midway,above]{$\mat{A}$};
  % output: collapsed onto x-axis
  \begin{scope}[xshift=6.0cm]
    \draw[->,gray!55] (-1.2,0)--(2.6,0) node[right]{};
    \draw[->,gray!55] (0,-1.7)--(0,1.9) node[above]{};
    % image is just the segment along x (sigma1=2)
    \draw[targetgreen,very thick] (-2,0)--(2,0);
    \draw[->,vecred,very thick] (0,0)--(2,0) node[below right,vecred]{$\sigma_1\vec{u}_1$};
    \fill[structurepurple] (0,0) circle (1.6pt);
    \node[structurepurple,below left] at (0,0) {$\vec{v}_2\mapsto\vec{0}$};
    \node at (0.7,-2.2) {image = a line};
  \end{scope}
  \draw[->,thick] (8.7,0)--(10.1,0) node[midway,above]{$\mat{A}^+$};
  % recovered input
  \begin{scope}[xshift=12.5cm]
    \draw[->,gray!55] (-1.7,0)--(1.9,0) node[right]{};
    \draw[->,gray!55] (0,-1.7)--(0,1.9) node[above]{};
    \draw[->,vecred,very thick] (0,0)--(1,0) node[below right,vecred]{$\vec{v}_1$};
    \fill[structurepurple] (0,0) circle (1.6pt);
    \node[structurepurple,above left,font=\small] at (0,0) {$\vec{0}$};
    \node[align=center] at (0.1,-2.2) {recovered:\\$\vec{v}_1$ restored, nothing in $\vec{v}_2$};
  \end{scope}
\end{tikzpicture}
\caption{The pseudoinverse on $\mat{A}=\protect\begin{bmatrix}2&0\\0&0\protect\end{bmatrix}$.
$\mat{A}$ stretches $\vec{v}_1$ by $\sigma_1=2$ and crushes $\vec{v}_2$ to zero,
so its image is just a line. $\mat{A}^+$ undoes the stretch on the reachable
direction (multiply by $1/\sigma_1$) and, having no way to recover the
crushed direction, returns nothing there. That is the least-squares,
minimum-norm answer.}
\label{ch29:fig-pseudoinverse}
\end{figure}

This closes a loop you opened chapters ago. The least-squares solution of an
overdetermined system, the projection onto a column space, the line of best fit --
all of them can be written as $\mat{A}^+\vec{b}$, and the SVD computes this solution stably even
when $\mat{A}\T\mat{A}$ is nearly singular (the case where the normal equations
quietly blow up). This is why serious numerical software solves least squares
through the SVD rather than by forming $\mat{A}\T\mat{A}$ directly.

\begin{example}[Pseudoinverse of a rank-deficient matrix]
\label{ch29:ex-pseudoinverse}
Let $\mat{A} = \begin{bmatrix} 2 & 0 \\ 0 & 0\end{bmatrix}$, which collapses the
plane onto the $x$-axis (Figure~\ref{ch29:fig-pseudoinverse}). It has no inverse.
Its SVD is trivial here:
$\mat{U} = \mat{V} = \mat{I}$, $\mat{\Sigma} = \diag(2,0)$. So
\[
  \mat{\Sigma}^+ = \begin{bmatrix} \tfrac12 & 0 \\ 0 & 0 \end{bmatrix},
  \qquad
  \mat{A}^+ = \begin{bmatrix} \tfrac12 & 0 \\ 0 & 0 \end{bmatrix}.
\]
Check the spirit of it: $\mat{A}^+$ undoes the doubling on the $x$-axis (multiply
by $\tfrac12$) and does the only reasonable thing with the crushed $y$-direction
-- leaves it at zero, since no input could ever have produced a nonzero
$y$-output.
\end{example}

\section{Symmetry, Eigenvalues, and the SVD}
\label{ch29:symmetric-link}

How does the SVD relate to the eigen-decomposition you already know? They coincide
in the nicest case and diverge instructively otherwise.

If $\mat{A}$ is \textbf{symmetric positive definite}, its spectral decomposition
$\mat{A} = \mat{Q}\mat{\Lambda}\mat{Q}\T$ \emph{is} its SVD: take $\mat{U} =
\mat{V} = \mat{Q}$ and $\sigma_i = \lambda_i$. The eigen-axes are already
orthonormal and the eigenvalues already positive, so there is nothing left to do.

If $\mat{A}$ is \textbf{symmetric but has a negative eigenvalue}, the SVD differs
by a sign. Suppose $\mat{A}\vec{q} = \lambda\vec{q}$ with $\lambda < 0$. The
singular value must be positive, $\sigma = \abs{\lambda} = -\lambda$, and to keep
$\mat{A}\vec{v} = \sigma\vec{u}$ intact we set $\vec{v} = \vec{q}$ but
$\vec{u} = -\vec{q}$ (or vice versa) so the sign lands in $\mat{U}$ instead of in
$\sigma$. In general: \textbf{the singular values of a symmetric matrix are the
absolute values of its eigenvalues}.

\begin{connectionbox}
\textbf{Polar decomposition.} Group the SVD as
\[
  \mat{A} = \underbrace{(\mat{U}\mat{V}\T)}_{\mat{Q}}
            \underbrace{(\mat{V}\mat{\Sigma}\mat{V}\T)}_{\mat{P}} = \mat{Q}\mat{P}.
\]
Then $\mat{Q} = \mat{U}\mat{V}\T$ is orthogonal (a pure rotation/reflection) and
$\mat{P} = \mat{V}\mat{\Sigma}\mat{V}\T$ is symmetric positive semidefinite (a
pure stretch along orthogonal axes). This is the matrix analogue of writing a
complex number as $z = e^{i\theta}\cdot r$: a rotation times a stretch. Engineers
use the polar decomposition to split a deformation into ``rigid rotation'' plus
``pure strain,'' and graphics programmers use it to extract the nearest true
rotation from a slightly corrupted one. It falls out of the SVD in a single line.
\end{connectionbox}

\begin{historybox}
The singular value decomposition was discovered independently at least four times, as mathematicians in different fields stumbled upon the same geometric truth from completely different directions.

\textbf{Who and when:} Eugenio Beltrami (1873) in Italy and Camille Jordan (1874) in France first derived the decomposition for real square matrices. James Joseph Sylvester (1889) in England rediscovered it independently fifteen years later, calling the singular values ``canonical multipliers.'' In 1907, Erhard Schmidt in Germany generalized the idea to integral operators in infinite dimensions. In 1936, the American psychometricians Carl Eckart and Gale Young proved the celebrated low-rank approximation theorem that made the SVD the mathematical foundation of data compression.

\textbf{What problem they were stuck on:} Beltrami and Jordan were studying \emph{bilinear forms}---functions that take two vector inputs and produce a single scalar, linear in each slot (such as formulas for kinetic energy, work, or geometric surface curvature). They asked: can one choose an independent orthogonal change of coordinates for the first vector and another for the second vector that diagonalizes the formula simultaneously? Sylvester was pursuing canonical forms in invariant theory, while Schmidt was solving Fredholm integral equations---the infinite-dimensional analogue of matrix equations, where discrete sums over vector coordinates become continuous integrals over functions. Eckart and Young were trying to extract a few fundamental psychological factors (such as verbal or spatial reasoning) from large tables of student test scores, leading them to ask: what is the single best rank-$k$ approximation to an empirical data matrix?

\textbf{The tools they had:} Pencil, paper, determinant expansions, and classical eigenvalue theory for quadratic forms. None of the early pioneers had access to computers. When electronic computers arrived in the 1950s, scientists immediately tried to compute the SVD using the textbook formula: form $\mat{A}\T\mat{A}$, find its eigenvalues, and take square roots. This triggered a severe numerical crisis. In finite-precision floating-point arithmetic, forming $\mat{A}\T\mat{A}$ squares the matrix's \emph{condition number} (the ratio of largest to smallest stretch, $\sigma_{\max}/\sigma_{\min}$, which measures sensitivity to perturbations). If a matrix has singular values spanning $1$ down to $10^{-8}$, $\mat{A}\T\mat{A}$ spreads them from $1$ down to $10^{-16}$---right at the edge of double precision, where small singular values that reveal near-singularity are wiped out by rounding noise.

\textbf{Controversy and resolution:} Sylvester's 1889 paper provoked priority friction, as Jordan had published the identical canonical form fifteen years earlier. Decades later, Eckart and Young's 1936 paper was found to have been largely anticipated by Schmidt's 1907 work on integral operators, leading many historians to refer to the low-rank approximation result as the Schmidt--Eckart--Young--Mirsky theorem. The computational bottleneck was finally resolved in 1965 by Gene Golub and William Kahan, who introduced a backward-stable algorithm that reduces $\mat{A}$ to \emph{bidiagonal form} (a matrix with nonzeros only on the main diagonal and the diagonal directly above it) using \emph{Householder reflections} (orthogonal mirrorings across hyperplanes) applied directly to $\mat{A}$, completely avoiding the dangerous $\mat{A}\T\mat{A}$.

\textbf{Impact:} The SVD applies to every matrix without restriction. Unlike the Jordan canonical form, which requires square matrices and breaks down when eigenvectors collide, the SVD exists unconditionally for any matrix, rectangular or defective, and provides an orthonormal coordinate frame for all four fundamental subspaces simultaneously. It gave science the pseudoinverse for ill-conditioned least squares, Principal Component Analysis for dimensionality reduction, and low-rank compression for modern deep learning.

\textbf{What is still open:} For two-dimensional matrices, the low-rank approximation problem is completely solved. But for \emph{tensors}---three-dimensional or higher data arrays, such as video volumes or multi-modal embeddings---the situation is fundamentally harder. Finding the best rank-$k$ approximation of a 3-way tensor is provably NP-hard, and the best rank-$k$ approximation can fail to exist because tensor rank is not topologically closed (a sequence of rank-2 tensors can converge to a rank-3 tensor). At the same time, computing \emph{sparse} singular vectors (forcing many coordinates to be exactly zero for interpretability in genetics and finance) is also NP-hard. Scalable, communication-efficient low-rank approximations for massive distributed datasets remain an active research frontier.
\end{historybox}

\section{How It Is Computed (and Why Not the Obvious Way)}
\label{ch29:computation}

The hand recipe -- diagonalize $\mat{A}\T\mat{A}$, take square roots -- is perfect
for a $2\times 2$ by hand and a disaster for a serious computation. The reason is
\emph{conditioning}: forming $\mat{A}\T\mat{A}$ squares the singular values, and
squaring spreads them out. A matrix whose singular values range from $1$ down to
$10^{-8}$ becomes, after $\mat{A}\T\mat{A}$, a matrix with eigenvalues from $1$
down to $10^{-16}$ -- right at the edge of what floating-point can distinguish from
zero. You can lose half your accuracy before you start.

\begin{trapbox}
Never compute an SVD (or a least-squares solution) in production by forming
$\mat{A}\T\mat{A}$ if you can avoid it. The squaring throws away precision exactly
where you need it most -- in the small singular values that tell you about
near-rank-deficiency. Real algorithms (Golub--Kahan bidiagonalization and its
descendants) work on $\mat{A}$ directly, never forming $\mat{A}\T\mat{A}$, and
keep full accuracy. When you call \texttt{numpy.linalg.svd} or LAPACK's
\texttt{dgesdd}, this is what is running.
\end{trapbox}

For a feel of the cost: a full SVD of an $m\times n$ matrix takes on the order of
$mn\cdot\min(m,n)$ arithmetic operations -- comparable to a few matrix
multiplications. For the giant, mostly-low-rank matrices of data science, one
rarely wants the full thing; \emph{randomized} and \emph{truncated} SVD
algorithms find just the top $k$ singular triples in roughly $O(mnk)$ time, which
is what makes SVD-based methods feasible on web-scale data.

\section{Summary: One Picture, Many Doors}
\label{ch29:summary}

\begin{focusbox}
Everything in this chapter hangs on one sentence and one equation.
\begin{itemize}
  \item \textbf{Geometry:} every matrix sends the unit sphere to an ellipsoid;
        $\mat{A} = \mat{U}\mat{\Sigma}\mat{V}\T$ is \emph{rotate, stretch,
        rotate}.
  \item \textbf{Pieces:} $\vec{v}_i$ = eigenvectors of $\mat{A}\T\mat{A}$;
        $\sigma_i = \sqrt{\lambda_i}$; $\vec{u}_i = \mat{A}\vec{v}_i/\sigma_i$.
  \item \textbf{Subspaces:} $\vec{v}$'s split into row space / null space;
        $\vec{u}$'s split into column space / left null space; $\sigma_i$ is the
        stretch from one to the other.
  \item \textbf{Approximation (Eckart--Young):} keep the top $k$ terms of
        $\sum \sigma_i \vec{u}_i\vec{v}_i\T$ for the best rank-$k$ matrix; error
        is the tail of the singular values.
  \item \textbf{Inversion:} $\mat{A}^+ = \mat{V}\mat{\Sigma}^+\mat{U}\T$ solves
        least squares, minimum norm, always.
\end{itemize}
\end{focusbox}

In the next chapter, we apply this factorization to empirical data. When you collect
multivariate observations -- gene expression levels, movie ratings, or financial
indicators -- the data points form an elongated cloud in high dimensions.
Center that cloud at its mean, stack the points as rows of a matrix, and the SVD
identifies the orthogonal axes along which the data varies most: the
\emph{principal components}. Principal Component Analysis is the direct
application of the SVD to sample covariance matrices. You already have the geometric foundation to see
exactly how it works.

\section{Exercises}
\label{ch29:exercises}

\begin{exercise}[Warmup: read off the SVD]
\label{ch29:ex-readoff}
Write the SVD $\mat{U}\mat{\Sigma}\mat{V}\T$ of each matrix by inspection (no
$\mat{A}\T\mat{A}$ needed):
(a) $\begin{bmatrix} 5 & 0 \\ 0 & 2\end{bmatrix}$;
(b) $\begin{bmatrix} 0 & 3 \\ 2 & 0\end{bmatrix}$ (hint: it is a scaling
composed with a swap of axes);
(c) the rotation $\begin{bmatrix} \cos\theta & -\sin\theta \\ \sin\theta & \cos\theta\end{bmatrix}$;
(d) Can any nonzero matrix have all of its singular values equal to zero? Why or why not?
What are the singular values in each case?
\end{exercise}

\begin{exercise}[A full $2\times 2$ by hand]
\label{ch29:ex-full2x2}
Compute the SVD of $\mat{A} = \begin{bmatrix} 1 & 1 \\ 0 & 1\end{bmatrix}$ (the
unit shear). Form $\mat{A}\T\mat{A}$, find its eigenvalues and eigenvectors, the
singular values, and the $\vec{u}$'s. Confirm that even though this matrix has
the single repeated eigenvalue $1$ and only one eigenvector, it has two distinct
singular values and a complete SVD. What does this say about the difference
between eigenvalues and singular values?
\end{exercise}

\begin{exercise}[Rank-one matrices]
\label{ch29:ex-rankone}
Let $\mat{A} = \vec{a}\vec{b}\T$ for nonzero vectors $\vec{a}\in\R^m$,
$\vec{b}\in\R^n$. Show that $\mat{A}$ has exactly one nonzero singular value, and
find it in terms of $\norm{\vec{a}}$ and $\norm{\vec{b}}$. Identify $\vec{u}_1$
and $\vec{v}_1$.
\end{exercise}

\begin{exercise}[Singular values vs.\ eigenvalues]
\label{ch29:ex-svd-vs-eig}
Give a $2\times 2$ matrix whose eigenvalues are both $1$ but whose larger
singular value is greater than $100$. (Hint: a strong shear.) Explain
geometrically how a matrix can fix the length of two directions yet stretch the
unit circle enormously.
\end{exercise}

\begin{exercise}[Determinant and singular values]
\label{ch29:ex-det}
Prove that for any square $\mat{A}$, $\abs{\det\mat{A}} = \prod_i \sigma_i$, using
$\mat{A} = \mat{U}\mat{\Sigma}\mat{V}\T$ and the facts that $\det$ is
multiplicative and $\abs{\det\mat{U}} = \abs{\det\mat{V}} = 1$. Deduce that
$\mat{A}$ is invertible if and only if all its singular values are positive.
\end{exercise}

\begin{exercise}[Frobenius norm and the singular values]
\label{ch29:ex-frob}
Using $\norm{\mat{A}}_F^2 = \tr(\mat{A}\T\mat{A})$, the SVD, and the fact that
trace is unchanged by orthogonal similarity, prove that
$\norm{\mat{A}}_F^2 = \sigma_1^2 + \cdots + \sigma_r^2$. Then explain why the best
rank-$k$ approximation leaves Frobenius error
$\sqrt{\sigma_{k+1}^2 + \cdots + \sigma_r^2}$.
\end{exercise}

\begin{exercise}[Spectral norm]
\label{ch29:ex-spectral}
The spectral norm is $\norm{\mat{A}}_2 = \max_{\norm{\vec{x}}=1}\norm{\mat{A}\vec{x}}$.
Using the SVD, prove that $\norm{\mat{A}}_2 = \sigma_1$, the largest singular
value, and that the maximizing $\vec{x}$ is $\vec{v}_1$. (Hint: write $\vec{x}$ in
the orthonormal basis $\vec{v}_1,\dots,\vec{v}_n$ and compute
$\norm{\mat{A}\vec{x}}^2$.)
\end{exercise}

\begin{exercise}[Pseudoinverse and least squares]
\label{ch29:ex-pinv}
For $\mat{A} = \begin{bmatrix} 1 & 0 \\ 0 & 0 \\ 0 & 0\end{bmatrix}$ and
$\vec{b} = \colvec{3\\4\\5}$, write the SVD of $\mat{A}$, form $\mat{A}^+$, and
compute $\hat{\vec{x}} = \mat{A}^+\vec{b}$. Verify directly that
$\hat{\vec{x}}$ minimizes $\norm{\mat{A}\vec{x} - \vec{b}}$ and explain
geometrically why the residual vector $\vec{b} - \mat{A}\hat{\vec{x}} = \colvec{0\\4\\5}$
lies entirely in the left null space of $\mat{A}$.
\end{exercise}

\begin{exercise}[Polar decomposition by hand]
\label{ch29:ex-polar}
Using the SVD found in Section~\ref{ch29:worked-2x2} for
$\mat{A} = \begin{bmatrix} 3 & 0 \\ 4 & 5\end{bmatrix}$, compute the polar
factors $\mat{Q} = \mat{U}\mat{V}\T$ and $\mat{P} = \mat{V}\mat{\Sigma}\mat{V}\T$.
Check that $\mat{Q}$ is orthogonal ($\mat{Q}\T\mat{Q} = \mat{I}$), that $\mat{P}$
is symmetric, and that $\mat{Q}\mat{P} = \mat{A}$.
\end{exercise}

\begin{exercise}[Why truncation is greedy-optimal]
\label{ch29:ex-greedy}
Argue, using the orthogonality of the $\vec{u}_i$ and $\vec{v}_i$, that the
``energy'' $\norm{\mat{A}}_F^2$ splits additively across the rank-one terms:
$\norm{\mat{A}}_F^2 = \sum_i \sigma_i^2$ with no cross terms. Explain in one or
two sentences why this additivity is exactly what makes ``keep the largest
$\sigma_i$'' the optimal way to spend a rank budget. (This is the intuition behind
Eckart--Young.)
\end{exercise}

\begin{exercise}[Challenge: SVD of $\mat{A}$ vs.\ $\mat{A}\T$]
\label{ch29:ex-transpose}
If $\mat{A} = \mat{U}\mat{\Sigma}\mat{V}\T$, write down the SVD of $\mat{A}\T$.
What happens to the singular values, the left singular vectors, and the right
singular vectors? Use this to explain in one sentence why $\mat{A}$ and
$\mat{A}\T$ always have the same singular values, the same rank, and why
$\mat{A}\T\mat{A}$ and $\mat{A}\mat{A}\T$ share their nonzero eigenvalues.
\end{exercise}
